\capitulo{2}{Objetivos del proyecto}

\section{Introducción}

El presente TFG tiene como objetivo principal el desarrollo de un \textbf{prototipo de laboratorio virtual de programación y robótica}  \textbf{enfocado en la enseñanza de la programación a través de un sistema de desafíos interactivos}  que permita a los estudiantes aprender programación mediante bloques y código escrito. Desarrollado con el motor de videojuegos \textbf{Unity 6} y el lenguaje de programación \textbf{C\#}, este laboratorio ofrece una plataforma  que permite a los estudiantes aprender a programar tanto mediante \textbf{código escrito} como con \textbf{bloques visuales}.

La educación en robótica y programación ha demostrado ser una herramienta clave en el aprendizaje del pensamiento computacional. Sin embargo, uno de los principales desafíos en la enseñanza de la robótica es la dependencia de hardware físico, lo que limita el acceso a estudiantes en contextos de escasos recursos y en zonas menos favorecidas.

Para resolver este problema, el laboratorio virtual ofrece una solución dual. Por un lado, los estudiantes pueden programar el robot en un entorno 3D utilizando un sistema de bloques de colores, similar a Scratch, que permite arrastrar y conectar instrucciones de forma intuitiva. Por otro lado, y aquí reside la principal innovación, también pueden resolver los mismos desafíos escribiendo comandos directamente en un lenguaje de programación textual simple, diseñado específicamente para ser fácil de aprender. Ambas formas de programar controlan al mismo robot virtual, inspirado en kits como LEGO WeDo 2.0, permitiendo a los usuarios elegir el método que prefieran o transicionar de uno a otro.

Para guiar el desarrollo de este prototipo, se establecieron los siguientes objetivos específicos:

\section{Objetivos funcionales}

Estos objetivos están relacionados con las funcionalidades que el prototipo de laboratorio virtual debe cumplir para ofrecer una experiencia educativa inicial.

\begin{enumerate}
	\item \textbf{Diseñar y desarrollar un sistema de programación para robótica educativa.}
	\begin{itemize}
		\item Crear un entorno de programación que combine la programación visual por bloques con la programación textual.
		\item Replicar el sistema de programación por bloques simulando a Scratch
		\item Desarrollar un Lenguaje de Dominio Específico (DSL) \cite{DSL} \cite{lenguaje_DSL}  textual para los comandos de bajo nivel del robot (mover, girar. repetir), y un \textbf{parser} para validar y procesar dichos comandos.
	\end{itemize}

\item \textbf{Integrar el sistema de programación con un entorno de simulación 3D.}
\begin{itemize}
	\item  Replicar la lógica fundamental de interacción de Scratch, permitiendo arrastrar, soltar y conectar bloques visuales en un área de trabajo.
	\item Desarrollar un motor de ejecución (CSharpRunner) que interprete la secuencia de bloques visuales y la traduzca en acciones para el robot simulado.
	\item Crear una librería de bloques básicos funcionales, incluyendo un bloque de evento para iniciar el programa, un bloque de control para repetir acciones y bloques de acción como mover pasos y girar
	\item Diseñar la gramática de un Lenguaje de Dominio Específico (DSL) simple y orientado a la robótica (mover, girar, repetir, ...).
	\item Desarrollar un parser que analice el código fuente del DSL, valide su sintaxis y lo transforme en una lista de instrucciones ejecutable.
	\item Crear una interfaz de desarrollo (IDE) específica para este modo, donde el usuario escribe comandos y recibe feedback visual inmediato de su programa.
\end{itemize}

\item \textbf{Implementar un prototipo de laboratorio virtual basado en desafíos.}
\begin{itemize}
	\item Crear un modelo 3D de un robot genérico inspirado en la estética de kits educativos como el robot Lego Wedo 2.0.
	\item Dotar al robot virtual de las funcionalidades de movimiento y rotación necesarias para completar los desafíos.
	\item Diseñar un sistema de gestión de desafíos que permita definir objetivos concretos, cargando dinámicamente las instrucciones y la configuración de la escena.
\end{itemize}
\end{enumerate}

\section{Objetivos no funcionales}

Estos objetivos se enfocan en las cualidades transversales de la solución, como su accesibilidad, rendimiento y usabilidad.

\begin{enumerate}
	\item \textbf{Accesibilidad y escalabilidad del prototipo.}
	\begin{itemize}
		\item Desarrollar una plataforma que pueda ejecutarse en dispositivos de hardware común y sin requerimientos técnicos elevados, con el objetivo de priorizar la portabilidad a plataformas Android (Tablets).
		\item Explorar estrategias de diseño y desarrollo que minimicen la dependencia de equipamiento costoso, haciendo la herramienta accesible a estudiantes en zonas con recursos limitados.

	\end{itemize}
	
	\item \textbf{Investigación y análisis.}
	\begin{itemize}
		\item Investigar y analizar plataformas existentes de programación y simulación robótica para comprender sus fortalezas y limitaciones.
		\item Definir qué aspectos de las soluciones existentes pueden ser mejorados o adaptados en la solución propuesta de este TFG, sentando las bases para futuras evoluciones del sistema.
	\end{itemize}
	
	\item \textbf{Adquisición de conocimientos técnicos.}
	\begin{itemize}
		\item Adquirir conocimientos y habilidades en el uso del motor Unity 6 y el lenguaje C\# para la creación de entornos interactivos y simulaciones educativas.
		\item Aplicar buenas prácticas en el diseño de software y en la integración de sistemas de programación visual dentro de un motor de simulación, incluyendo la aplicación de patrones de diseño (MVC, Observador, Fábrica) y la gestión eficiente de estructuras de datos.
	\end{itemize}
\end{enumerate}
